\summarychapter{5}{Arguing About the Event}{The Rise of Orthodoxy}

\begin{inbrief}
Catholic Christianity was orthodox as well as universal. The chapter explains what theology, heresy
and orthodoxy are, shows how the apostles already drew lines against false gospels, and describes the
church's first great struggle over ideas, the conflict with Gnosticism, out of which came the
Apostles' Creed.
\question{What did the first Christians hold to be true, and why would they accept no other version of the faith as orthodox?}
\end{inbrief}

Gandhi said he would not care if the man Jesus had never lived, since for him the truth of the Sermon on the Mount did not depend on it. Many professing Christians think the same way: they separate what Jesus said
from who he was and prefer behavior to belief. The early believers called that a betrayal. The gospel
was good news about the Event, so belief came first and behavior followed, and belief in who Jesus was
became the test of true Christianity. Persecution was the outward enemy; a subtler danger, ``another
gospel,'' came from within.

\ssection{Faith and Theology}
Many dislike theology (``I love flowers but I hate botany''), often with reason, since it can be dull
or ruthless. But the answer to bad theology is good theology, not none. Theology (\emph{theos}, God;
\emph{logos}, rational thought) is thinking about religion; religion is belief in God and the effort
to live by it. Wrong thinking is heresy, which is not necessarily bad religion but may lead to it.
Heretics served the church unintentionally: their attempts forced it to produce good theology, a
rounded statement of biblical revelation. That theology, held by the overwhelming majority and
expressed in the creeds, is called orthodox. Doctrine was hammered out by people on the church's
``work crew,'' usually because a heresy threatened the central faith. Hence orthodoxy
\emph{developed}: revelation was given once for all, but theology is the human understanding of it,
in the language and thought of each time and culture. One complication was the meeting of Hebrew
thought, rooted in covenant and history, with Hellenistic reasoning. The church did not make its
beliefs depend on logical demonstration; it believed because God had revealed himself. The first
Christians, all Jews, proclaimed Jesus as the Messiah, stressed the resurrection as God's verdict on
the crucified, and described him in Old Testament images: Passover Lamb, second Adam, Son of David,
the rejected stone become the cornerstone.

\begin{sidebox}{Orthodoxy and Heresy}
``Orthodox'' can name a branch of Christianity, right belief in general, or belief officially adopted
by the church; a heretic rejects the last, as Arius rejected Nicaea's teaching that Jesus is fully
divine. Official doctrine varied, and decisions such as Nicaea's were disputed in their day. Still,
there was an early \emph{functional} orthodoxy held by consensus, visible in the ``rule of faith,'' a
short statement of the Christian story. It is true that no one can break with a doctrine before it is
defined, but wrong to conclude that there were no heretics before the councils: the church had
already rejected Marcion, who cut Jesus off from the Old Testament God, and modalism, which made
Father, Son and Spirit mere roles of God.
\end{sidebox}

\ssection{False Gospels}
The apostles condemned rival gospels outright. Galatians curses those who add Jewish legal
requirements; 1 John requires belief that Christ came ``in the flesh''; 1 Corinthians makes the
historical resurrection the basis of salvation. Believers met the central truths in baptism, soon
given in the name of the Trinity (by Justin's time, with questions about Father, Son and Spirit), in
short summaries of apostolic teaching and in hymns such as 1 Timothy 3:16. Because the worship of
Jesus was central, the decisive line ran through the doctrine of Christ. John's gospel fights on two
fronts: against those who doubted that Jesus was fully God, and against those who denied his real
flesh and blood. Both errors are known from other sources. The \emph{Ebionites}, a Jewish-Christian
sect, held that Jesus was a mere man who became Messiah by his obedience to the Law.
\emph{Docetism} (from the Greek ``to seem'') held that Christ was a phantom who only seemed to
suffer. God in the flesh has always seemed nonsense, and people keep trying to turn the story into a
timeless myth.

\ssection{Men Who Knew}
The most ambitious attempt was Gnosticism, an umbrella term for movements each offering
enlightenment through a teacher who possessed the \emph{gnosis}, the knowledge of the way of life.
Feelings ran high: Polycarp told how the apostle John fled the baths at Ephesus on seeing the Gnostic
Cerinthus inside. Gnostics shared with the church the ideas of salvation, a supreme deity and
heavenly beings, which is why they hovered at its edges and entered it, hoping to purify these ideas
of a ``materialistic'' Christianity. Their basic belief was dualism: good and evil are two cosmic
forces, and evil is matter. A creator must therefore be wicked. The supreme being is far above
creation; from him came a chain of emanations, each weaker than the last, ending in one powerful
enough to create and foolish enough not to see that creation is wrong---the God of the Jews. Light
from heaven had become trapped in matter and needed redeeming. The Gnostic gospels show a Jesus who
is one of the lesser deities, cut off from his Jewish setting, without miracles or resurrection,
passing on knowledge by sayings. Readers today often find in him a tolerant, inclusive figure like
themselves, despite the saying against women that ends the \emph{Gospel of Thomas}. The church knew
these writings and rejected them: they may hold genuine sayings, but they give the wrong overall
impression of Jesus. All Gnostics agreed that the Christ could not be human, so the first great test
of faith in the Event was a denial not of Christ's deity but of his humanity. They also taught a kind
of predestination: an upper class living by knowledge, a lower class by faith, and a third incapable
of gnosis.

\begin{sidebox}{What Gnostics Believe}
The origins of Gnosticism are disputed, but its second-century teachers, refuted by Irenaeus,
Tertullian and Origen, shared five ideas: (1) a moral and metaphysical dualism, spirit pure and matter
evil; (2) a supreme god who is not the creator, the material world being the work of disorder or of a lesser divinity (called the demiurge and identified with the God of the Old Testament or with Jesus); (3) a divine element in humanity, or in
the chosen, which is the true self, an alien in the body; (4) a series of mediators in successive
heavens; (5) salvation as the spirit's liberation from matter through knowledge.
\end{sidebox}

\begin{sidebox}{The Gospel of Thomas}
Gnostic accounts of Jesus are late, mostly from the late second to the fifth century. The possible
exception is the \emph{Gospel of Thomas}, found in Egypt in 1945 in a fourth-century Coptic copy, with
Greek fragments from about 200 onward. It was probably written in Syriac. Some date it early, but its
use of Syriac traditions points to 150--200, most likely 175--200. An early date would change little:
the church was actively guarding the genuine memory of Jesus.
\end{sidebox}

\begin{sidebox}{Politics and Interpretation}
Two strategies of interpretation met in the \emph{Gospel of Thomas}. One assumes that history's
winners, here a bishop-led Christianity, wrote the losers out of the record. The other suspects that
a counterfeit Christianity replaced the genuine one; Thomas Jefferson thought the four gospels a
corrupted version of Jesus. Together they produce a ``suppressed'' Gnostic Christianity.
\emph{The Da Vinci Code} is a popular, unscholarly version, which implausibly dates the turn to
traditional Christianity to the fourth century.
\end{sidebox}

\ssection{The Dangers of Knowing}
Gnosticism is a warning to all who would free the gospel from ``outmoded'' Jewish ideas of God and
history or lift it from faith to a higher knowledge to impress important people. The Gnostic made Christ fit the science and philosophy then in fashion, and in doing so gave up the Event and with it the gospel; nothing fades
faster than the latest theories of the enlightened. Had Gnosticism won, the heritage from Judaism
would have been surrendered, the message for all would have shrunk to talk among a few, and Christ
would have been lost among the gods of the mystery religions. Gnostics were hard to refute because
they claimed secret teaching from Jesus or special revelations.

The church's answer clarified its own convictions, best summed up in the Apostles' Creed. It was not
written by the apostles but grew from a second-century baptismal confession at Rome, the Old Roman
Creed. Built on the Trinity without explaining it, the Creed is concerned with how God relates to the
world. ``God Almighty'' (later, ``Maker of heaven and earth'') rules out the view that matter is bad or was made by a bad god; as C.~S.~Lewis said, God likes matter: he invented it. ``Born of the
Virgin Mary,'' ``crucified under Pontius Pilate'' and ``buried'' affirm the full humanity of
Jesus---the Gnostic stumbled at ``born,'' as moderns stumble at ``Virgin.'' Redemption comes by
God's action in history, not by secret knowledge: not gnosis but Event. ``The resurrection of the
flesh'' insists that a person is a whole, not a good soul in an evil body. People need salvation not
because they are imprisoned in bodies but because they willfully choose their own way; the evil is in
the affections. That is why, for catholic Christianity, Gandhi is wrong: humanity needs not a teacher
but a Savior. (The chapter also carries a chart, ``An Overview of Gnosticism.'')

\begin{sidebox}{Answers to Gnosticism}
The ancient church met Gnostic and similar challenges with clergy, creed and canon. \emph{Clergy}:
its leaders stood in direct connection with the men Jesus had called and taught, who passed on the
message and left a written witness. \emph{Creed}: the story was condensed early into simple summaries;
those of Tertullian, Irenaeus and Origen, from different regions, agree remarkably---the Father sent
the Son to reclaim the world through his death and resurrection and the continuing work of the Spirit.
\emph{Canon}: the four Gospels, firmly held by the main body of the church within the second century,
give the right overall impression of Jesus.
\end{sidebox}

\begin{sidebox}{Understanding Gnosticism Today}
Two comparisons help. The church's gospels were written thirty to sixty-five years after Jesus, as
close as a professor of fifty-five is to the Vietnam or Korean wars, within living memory; the
earliest Gnostic gospel stands about 140 years away, as far as that professor is from the Civil War.
And reconstruction matters: books on Lincoln differ on details but share an outline, while a fantasy
about Lincoln the vampire hunter replaces and supplements the old story to make a quite different
one. The Gnostic versions of Jesus do the same.
\end{sidebox}

\begin{keyfacts}
\keypeople{Gandhi; the apostle John; Cerinthus; Polycarp; Marcion; Arius; Justin; Irenaeus;
Tertullian; Origen; Thomas Jefferson; C.~S.~Lewis; Charles Bigg.}
\keydates{%
  \kd{c. 150--200}{Probable date of the \emph{Gospel of Thomas}}
  \kd{2nd century}{The Old Roman Creed, ancestor of the Apostles' Creed, used at baptisms in Rome}
  \kd{1945}{The \emph{Gospel of Thomas} discovered in Egypt}}
\keyterms{Theology; heresy; orthodoxy; rule of faith; Ebionites; Docetism; Gnosticism, \emph{gnosis};
dualism; emanations; demiurge; modalism; Apostles' Creed; clergy, creed and canon.}
\end{keyfacts}
